\section{6차시 --- 랜덤 함수와 가위바위보 게임}
\label{sec:ch06}

본 차시는 "난수"라는 새로운 개념을 도입한다. 학습자는
\texttt{random.randint()}으로 매번 다른 결과를 만들어내는 비결정적
코드를 작성하며, 이를 게임 시나리오(가위바위보 / 3판 2선승제 / 타임어택
/ 표정 변화)로 발전시킨다.

\subsection{미션 1 --- 랜덤 가위바위보 시작!}
\label{subsec:ch06-m1}
\missionmeta{LED 3개}{RANDOM}

\begin{storybox}
\sayalbi{삐릿! 아레스, 심심한데 저와 가위바위보 게임 어때요?}
\sayares{오, 재미있을 것 같아! 그럼 알비는 LED 불빛으로, 나는 손으로 가위바위보 하는 거지?}
\end{storybox}

\subsubsection*{학습 목표}
\begin{itemize}[leftmargin=1.4em]
  \item LED 3개에 가위/바위/보 스티커 부착
  \item 랜덤으로 1$\sim$3 중 숫자 하나 선택
  \item 선택된 번호 LED 1개만 켜기
\end{itemize}

\begin{labbox}{샘플 코드 --- 가위바위보 랜덤}
\begin{minted}{python}
# 가위바위보 랜덤
import random
choice = random.randint(1, 3)
led_on(choice, 1.0)
time.sleep(2)
led_off(choice)
\end{minted}
\end{labbox}

\subsection{미션 2 --- 가위바위보 하나 빼기!}
\label{subsec:ch06-m2}
\missionmeta{LED 3개}{RANDOM}

\begin{storybox}
\sayalbi{삐릿♪ 아레스, 하나 빼기 게임 알아요? 가위바위보, 하나 빼기$\sim$ 하면서 노는 거예요!}
\sayares{오, 재밌겠다! LED 3개가 다 켜졌다가 2개 남고, 그다음 1개만 남는 거지?}
\end{storybox}

\subsubsection*{학습 목표}
\begin{itemize}[leftmargin=1.4em]
  \item LED 3개 동시에 3번 깜빡임 (시작 신호)
  \item \texttt{one\_out} 블록으로 2개 랜덤 켜기
  \item 1.5초 후 랜덤으로 1개 끄기 $\rightarrow$ 1개만 남음!
  \item ``가위바위보, 하나 빼기$\sim$!'' 게임 완성
\end{itemize}

\begin{labbox}{샘플 코드 --- 하나 빼기}
\begin{minted}{python}
# 시작 신호 (3번 깜빡)
for i in range(3):
    led_on(1, 1.0); led_on(2, 1.0); led_on(3, 1.0)
    time.sleep(0.3)
    led_off(1); led_off(2); led_off(3)
    time.sleep(0.3)
# 하나 빼기!
one_out()   # 랜덤 2개 켜기 → 1.5초 후 1개 끄기
\end{minted}
\end{labbox}

\subsection{미션 3 --- 알비를 이겨라!}
\label{subsec:ch06-m3}
\missionmeta{LED 3개}{RANDOM}

\begin{storybox}
\sayares{알비, 이번엔 진짜 게임이야! 알비가 5번 내면, 나는 매번 이기는 손을 빠르게 내는 거야!}
\sayalbi{삐릿- 도전 받아드려요! 시작 신호 3번 깜빡이고, 알비가 2초마다 랜덤으로 5번 낼게요!}
\end{storybox}

\subsubsection*{학습 목표}
\begin{itemize}[leftmargin=1.4em]
  \item 시작 신호: LED 3개 3번 깜빡임
  \item 알비가 2초마다 랜덤으로 5번 손을 냄
  \item 아레스는 이기는 손을 빠르게 내야 해요!
  \item 종료 신호: LED 3개 3번 깜빡임
\end{itemize}

\begin{labbox}{샘플 코드 --- 알비를 이겨라}
\begin{minted}{python}
# 시작 신호
for i in range(3):
    led_on(1,1); led_on(2,1); led_on(3,1)
    time.sleep(0.2)
    led_off(1); led_off(2); led_off(3)
    time.sleep(0.2)
# 알비 5번 랜덤 선택
for i in range(5):
    random_led()
# 종료 신호
for i in range(3):
    led_on(1,1); led_on(2,1); led_on(3,1)
    time.sleep(0.2)
    led_off(1); led_off(2); led_off(3)
    time.sleep(0.2)
\end{minted}
\end{labbox}

\subsection{미션 4 --- 가위바위보 왕중왕!}
\label{subsec:ch06-m4}
\missionmeta{LED 3개}{RANDOM}

\begin{storybox}
\sayares{알비, 드디어 마지막 결전이야! 3라운드를 해서 모두 이기면 화성 가위바위보 챔피언!}
\sayalbi{삐릿♪ 매 라운드 3번 깜빡이면 알비가 손을 낼게요! 이기는 손을 빠르게 내세요! 도전!}
\end{storybox}

\subsubsection*{학습 목표}
\begin{itemize}[leftmargin=1.4em]
  \item 3라운드 반복 구조 완성
  \item 매 라운드: 가위바위보$\sim$ 신호 3번 $\rightarrow$ 알비 랜덤 선택
  \item 이기는 손을 빠르게 내세요!
  \item 마지막엔 축하 댄스: LED 5번 빠르게 깜빡
\end{itemize}

\begin{labbox}{샘플 코드 --- 가위바위보 왕중왕}
\begin{minted}{python}
# 3라운드 가위바위보 왕중왕!
for round in range(3):
    # 준비 신호: 가위바위보~
    for i in range(3):
        led_on(1,1); led_on(2,1); led_on(3,1)
        time.sleep(0.3)
        led_off(1); led_off(2); led_off(3)
        time.sleep(0.3)
    random_led()     # 알비의 선택
    time.sleep(0.3)
# 챔피언 축하 댄스!
for i in range(5):
    led_on(1,1); led_on(2,1); led_on(3,1)
    time.sleep(0.15)
    led_off(1); led_off(2); led_off(3)
    time.sleep(0.15)
\end{minted}
\end{labbox}

\begin{keypoint}
난수는 "결정적 코드도 다른 결과를 만들 수 있다"는 점에서 학습자에게
강한 인상을 남긴다. \texttt{random.randint(a,b)}의 의미와, 동일 코드를
실행해도 매번 결과가 달라지는 현상을 직접 관찰하게 하는 것이 본 차시의
핵심 학습 포인트이다.
\end{keypoint}
